%
% Introduction
%

% !TEX root = ../../main.tex

\chapter{Einleitung}

  In der Technik, insbesondere in der Informatik und Robotik,
  ist das Lösen und Optimieren von Problemstellungen eine aufwendige Planungsaufgabe.
  Das Lösen dieser Problemstellungen erfordert viel Fachwissen, kostet Zeit und Ressourcen.
  Sobald der Lösungsraum eines Problems zu komplex wird, stossen die formalen Lösungsmethoden an ihre Grenzen.
  Oft spielt dabei die Anzahl der Freiheitsgrade eine entscheidende Rolle.
  \\
  Alternativ kann eine Lösung durch evolutionäre Algorithmen gesucht werden.
  Die Algorithmen produzieren Populationen, welche die möglichen Lösungen beinhalten.
  Das Konzept ist der natürlichen Evolution nachempfunden, wo durch ``Trial and Error'' neue Lösungen generiert werden.
  Mit dem Prinzip ``surival of the fittest'' werden immer bessere Näherungen für ein Problem produziert.
  \\
  Die Kombination aus Reproduktion, Vererbung, Variation und Selektion
  übernimmt dabei die Rolle des Lösens und Optimierens einer Problemstellung.
  \\
  Dawkins~\cite{book:universalDarwinism} definiert die vier Begrifflichkeiten~(auch universeller Darwinismus genannt) wie folgt:

  \begin{itemize}

    \item Reproduktion/Vererbung: Die Individuen müssen fähig sein sich fortzupflanzen.
      Der Nachwuchs welcher erzeugt wird, muss ebenfalls fortpflanzungsfähig sein und erbt Eigenschaften der Eltern.

    \item Variation: Es muss ein Mechanismus existieren~(Mutation~\vref{sub:mutFunction}),
      welcher neue Variationen in die Nachkommen einführen kann.

    \item Selektion: Die vererbten und variierten Merkmale müssen die Reproduktionsfähigkeit beeinflussen.
      Durch Selektion werden die am besten geeigneten Individuen selektiert.

  \end{itemize}

  In dieser Arbeit wird versucht mit evolutionären Algorithmen Individuen zu entwickeln,
  welche das Gehene lernen und einen Parcours bewältigen können.

  \section{Ausgangslage}

    \todo[inline]{Nennt bestehende Arbeiten/Literatur zum Thema -> Literaturrecherche}
    \todo[inline]{Stand der Technik: Bisherige Lösungen des Problems und deren Grenzen}

    % bestehende arbeiten
    % robotik -> nasa, steuerung roboter, informatik
    % stand der technik, aktueller ansatz
    % literatur erwähnen
    % Altes Problem der Robotik.
    % TODO moars

    Novel Approach to Content Based Image Retrieval Using Evolutionary Computing~\cite{Imran2014}.
    \\
    Evolutionary robotics: the biology, intelligence, and technology of self-organizing machines~\cite{Nolfi2001}.

    %
    %
    %
    %

    Evolutionäre Algorithmen sind bereits in vielen Bereichen im Einsatz.
    Besonders bei Problemstellungen wo eine approximative Lösung ausreichend ist,
    wird dieses Verfahren eingesetzt.
    \\
    Die NASA verwendet evolutionäre Algorithmen um Antennen zu entwerfen,
    die im Weltraum zur Funkkommunikation verwendet werden~\cite{Hornby2006}.
    Dabei wurde die Antenne auf gute Sendeleistung bei geringer Energieaufnahme optimiert.
    \\
    Als theoretische Grundlagen wurden die Bücher ``Bio inspired artificial intelligence''~\cite{book:bioInspired} und
    ``Evolutionäre Algorithmen''~\cite{book:evAlgo} verwendet.

  \section{Aufgabenstellung}

    ``Ziel der BA ist es,
    in einer virtuellen Umwelt mit Hilfe evolutionärer Algorithmen künstliche Tiere entstehen zu lassen.
    Evolviert werden die Geometrie und die Steuerung der Tiere,
    d.h.\ die Anordnung der Beine und deren Kontrolle.
    Als Mass für die Fitness wird eine Funktion der Zeit, welche das Tier braucht,
    um einen gegebenen Parcours zurückzulegen, verwendet.
    Als Randbedingungen vorgegeben sind: erstens die Forderung einer physikalisch sinnvollen Bewegung,
    zweitens eine obere Grenze für die total aufgewandte Energie und drittens eine Beschränkung der abgegebenen Leistung.
    Damit die Forderung nach einer physikalischen Bewegung erfüllt werden kann,
    wird die Bewegung mit Hilfe einer Physik-Engine simuliert.
    Diese Engine wird von den Studierenden ausgewählt und kann als Blackbox verwendet werden.
    \\
    Die BA beinhaltet folgende Forschungsfragen: 1) Wie kann eine Steuerung der Bewegung implementiert werden?
    2) Wie kann diese Steuerung evolviert werden? 3) Wie kann die Geometrie der Tiere evolviert werden?
    \\
    Erfüllungskriterien: A) Die Steuerung von Tieren mit vorgegebener Geometrie (6 Beine) wird evolviert,
    sodass ein Parcours mit vorgegebenen Eigenschaften bewältigt werden kann.
    B) In einem zweiten Schritt wird sowohl die Geometrie als auch die Steuerung evolutiv weiterentwickelt;
    dies mit dem Ziel, die für einen Parcours benötigte Zeit zu reduzieren.''~\cite{Task}

  \section{Zielsetzung}

    Das Ziel dieser Arbeit ist sechsbeinigen künstlichen Tieren mit Hilfe eines evolutionären Algorithmus,
    in einer zweidimensionalen Umgebung das Gehen beizubringen.
    Es soll nicht nur die Steuerung der Beine evolviert werden~(Gehen lernen),
    sondern auch die Form des Körpers und die Abmessung der Beine.
    \\
    Folgende Fragestellungen werden untersucht:

    \begin{enumerate}
      \item Wie kann eine Steuerung der Bewegung implementiert werden?\label{item:frageSteuerungBeine}
      \item Wie kann diese Steuerung evolviert werden?
      \item Wie kann die Geometrie der Tiere evolviert werden?
      \item Wie sieht der Bewegungsablauf und die Geometrie eines evolvierten Tieres aus?
      \item Nimmt die Diversität~(\vref{sub:diversity}) mit zunehmenden Generationen nur ab?
      \item Unterscheidet sich ein Individuum, welches mit dem Ziel der allgemeinen Lösung gefunden wurde,
            von einem das mit dem Evolvieren von Evolvierbarkeit gefunden wurde (\vref{sec:Konfiguration})?
    \end{enumerate}

    Eine bestehende \gls{PhysicsEngine} darf für das Simulieren der Tiere in einer virtuellen Umgebung eingesetzt werden.
    Der Hauptfokus liegt auf dem Evolvieren der artifiziellen Tiere und
    nicht auf dem Ausbau der verwendeten \gls{PhysicsEngine}.
    Ein Werkzeug in Form einer selbst programmierten Applikation soll erstellt werden,
    um den evolutionären Algorithmus zu implementieren.
    \\
    Weitere Hypothesen~(\vref{subsub:hypoKp} und~\vref{subsub:hypoMut})
    werden im Kapitel \nameref{sec:method}~(\vref{sec:method}) erwähnt,
    da die erforderlichen Begrifflichkeiten dann eingeführt sind.

  \section{Anforderungen}

    \subsection{Fitnessfunktion}

      Die Fitness ist ein Mass für den Grad der Anpassung eines Individuums an die Umgebung.
      In frühen Generationen können Tiere existieren, die sich nicht fortbewegen können.
      Die Zeit, welche ein Individuum für die Bewältigung eines Parcours (\vref{subsub:parcours}) brauchen würde,
      wäre somit unendlich.
      Um dieses Problem zu umgehen, wird als Mass für die Fitness die zurückgelegte Strecke verwendet.
      Für die Bewältigung des Parcours steht dem Individuum eine fixe Anzahl Zeiteinheiten zur Verfügung.

    \subsection{Parcours\label{subsub:parcours}}

      Die Individuen müssen während der Simulation eine Strecke zurücklegen.
      Diese Strecke und ihre Begebenheiten werden als Parcours bezeichnet.
      Die~\vref{fig:scetchParcours} zeigt eine Skizze des Parcours.
      Mit zunehmenden Generationen soll die Schwierigkeit des Parcours steigen.
      Die Schwierigkeit eines Parcours wird durch die maximale Steigung definiert.
      Der Parcours soll zufällig generiert werden.

      \begin{figure}[H]
        %
% Concept parcours
%

% !TEX root = ../main.tex

\begin{tikzpicture}[scale=\textwidth/20cm]

  \coordinate (P0) at (0,0);
  \coordinate (P1) at (0,1);
  \coordinate (P2) at (4,1);
  \coordinate (P3) at (7,2);
  \coordinate (P4) at (9,2);
  \coordinate (P5) at (10,1.5);
  \coordinate (P6) at (12,3);
  \coordinate (P7) at (14,3);
  \coordinate (P8) at (16,2.2);
  \coordinate (P9) at (20,2.2);
  \coordinate (P10) at (20,0);

  % Draw parcours
  \draw [fill=black] (P0) -- (P1) -- (P2) -- (P3) -- (P4) -- (P5) -- (P6) -- (P7) -- (P8) -- (P9) -- (P10) -- cycle;

\end{tikzpicture}

        \caption{Skizze Parcours\label{fig:scetchParcours}}
      \end{figure}

    \subsection{Limiten\label{sub:IntroReqLimit}}

      Limiten verhindern, dass der Parcours durch triviale Lösungen wie z.B. unendlich grosse Beine bewältigt werden kann.
      Ebenfalls schränken sie den Lösungsraum ein und helfen gewinnbringende Lösungen zu finden.
      \\
      Für folgende Komponenten werden Limiten eingeführt:

      \begin{itemize}
        \item Winkelgeschwindigkeit auf den Gelenken: Verhindert Beschleunigung ins unendliche.
        \item Beinhöhe: Verhindert die triviale Überwindung von Hindernissen.
        \item Beinbreite: Verhindert unbrauchbare Lösungen
        \item Masse des gesamten Individuums: Vereinfachung der Problemstellung
        \item Fläche des Körpers: Verhindert unendlich grosse Individuen
      \end{itemize}
